\paragraph{Notación.} Se usan las bases
\[
  \{\ket0,\ket1\}\qquad\text{y}\qquad\{\ket+,\ket-\},\qquad
  \ket\pm:=\frac{\ket0\pm\ket1}{\sqrt2}=\frac{e^{it}\pm e^{2it}}{\sqrt2}.
\]
El estado del sistema es
\[
  \Psi(t,s)=\frac1{\sqrt2}\bigl(e^{i(t+s)}+e^{2i(t+s)}\bigr)=\frac1{\sqrt2}\bigl(\ket0\otimes\ket0+\ket1\otimes\ket1\bigr).
\]
Según el enunciado, si Ana mide en una base ortonormal $\{u_0,u_1\}$ de $H_1$, el resultado $k$ tiene probabilidad $p_k=\|\Psi_k\|^2$, donde $\Psi_k(s)=\bk{u_k}{\Psi(\cdot,s)}$, y la onda de Bruno pasa a ser $\varphi_k=\Psi_k/\sqrt{p_k}$.

\paragraph{Medición de Ana en la base $\{\ket0,\ket1\}$.} Como $\bk{\ket0}{\Psi(\cdot,s)}=\frac1{\sqrt2}\ket0(s)$ y $\bk{\ket1}{\Psi(\cdot,s)}=\frac1{\sqrt2}\ket1(s)$, se obtienen $p_0=p_1=\frac12$, con $\varphi_0=\ket0$ y $\varphi_1=\ket1$.

\paragraph{Medición de Ana en la base $\{\ket+,\ket-\}$.} Como $\bk{\ket\pm}{\ket0}=\frac1{\sqrt2}$ y $\bk{\ket\pm}{\ket1}=\pm\frac1{\sqrt2}$,
\[
  \Psi_\pm(s)=\frac1{\sqrt2}\Bigl(\frac1{\sqrt2}\ket0(s)\pm\frac1{\sqrt2}\ket1(s)\Bigr)=\frac12\bigl(\ket0\pm\ket1\bigr)(s).
\]
Entonces $\|\Psi_\pm\|^2=\frac14+\frac14=\frac12$, es decir, $p_\pm=\frac12$, y $\varphi_\pm=\frac{\ket0\pm\ket1}{\sqrt2}=\ket\pm$.

Los resultados de Ana se resumen en la tabla siguiente.
\begin{center}
\renewcommand{\arraystretch}{1.3}
\begin{tabular}{cccc}
  Base de Ana & Resultado $k$ & $p_k$ & Estado de Bruno $\varphi_k$\\\hline
  $\{\ket0,\ket1\}$ & $0$ & $\frac12$ & $\ket0$\\
  $\{\ket0,\ket1\}$ & $1$ & $\frac12$ & $\ket1$\\
  $\{\ket+,\ket-\}$ & $+$ & $\frac12$ & $\ket+$\\
  $\{\ket+,\ket-\}$ & $-$ & $\frac12$ & $\ket-$
\end{tabular}
\end{center}

\paragraph{Probabilidades de Bruno.} Sea $\{w_0,w_1\}$ cualquier base ortonormal de $H_1$, en la que Bruno mide. Como $\{\ket0,\ket1\}$ es base ortonormal, para cualquier vector unitario $w$ se cumple $|\bk{w}{\ket0}|^2+|\bk{w}{\ket1}|^2=\|w\|^2=1$; lo mismo ocurre con la base $\{\ket+,\ket-\}$. Por tanto, si Ana mide en la base $\{\ket0,\ket1\}$,
\[
  P(\text{Bruno obtiene } w_j)=\tfrac12\,|\bk{w_j}{\ket0}|^2+\tfrac12\,|\bk{w_j}{\ket1}|^2=\tfrac12 ,
\]
y si Ana mide en la base $\{\ket+,\ket-\}$, la probabilidad también es
\[
  \tfrac12\,|\bk{w_j}{\ket+}|^2+\tfrac12\,|\bk{w_j}{\ket-}|^2=\tfrac12 .
\]
Bruno ve, entonces, probabilidades iguales a $\frac12$ para cualquier base suya y sin importar la base de Ana. La respuesta a la pregunta es que no dependen.

\paragraph{Reescritura con kets.} Con kets el cálculo muestra la razón. Como Bruno no conoce el resultado de Ana, su estado promedio es
\[
  \sum_k p_k\,\ket{\varphi_k}\bra{\varphi_k}.
\]
Para la base $\{\ket0,\ket1\}$ esta suma es $\frac12\ket0\bra0+\frac12\ket1\bra1=\frac12 I$, y para la base $\{\ket+,\ket-\}$ es $\frac12\ket+\bra++\frac12\ket-\bra-=\frac12 I$. Las dos sumas dan el mismo operador. Como la probabilidad de cada resultado $w_j$ de Bruno es $\bra{w_j}\bigl(\tfrac12 I\bigr)\ket{w_j}=\tfrac12$, Bruno no puede saber qué base eligió Ana. Lo que sí depende de la base de Ana son las probabilidades conjuntas, es decir, cómo se relacionan los resultados de los dos, y eso solo se ve cuando se comparan ambos resultados.

\paragraph{Implementación en Qiskit.} El circuito usa dos qubits: el qubit $0$ representa la variable $t$ (Ana) y el qubit $1$ la variable $s$ (Bruno). Como $\ket0$ corresponde a $e^{it}$ y $\ket1$ a $e^{2it}$, una compuerta de Hadamard sobre el qubit $0$ seguida de una CNOT prepara $\Psi$ a partir de $\ket{00}$. Si Ana mide en la base $\{\ket+,\ket-\}$, se aplica una Hadamard al qubit $0$ antes de medir en la base computacional; lo mismo hace Bruno con el qubit $1$. Como las operaciones de Ana actúan sobre un qubit y las de Bruno sobre el otro, medir ambos al final da la misma estadística que medir primero a Ana y después a Bruno. El código es el siguiente.
\begin{verbatim}
import numpy as np
from qiskit import QuantumCircuit
from qiskit.primitives import StatevectorSampler

SHOTS = 4000
# Un solo generador: cada experimento toma muestras nuevas
rng = np.random.default_rng(2026)
sampler = StatevectorSampler(seed=rng)

def circuito(base_ana, base_bruno):
    # Qubit 0 = variable t (Ana), qubit 1 = variable s (Bruno)
    qc = QuantumCircuit(2, 2)
    qc.h(0)                  # (|00> + |10>)/sqrt(2)
    qc.cx(0, 1)              # Psi = (|00> + |11>)/sqrt(2)
    if base_ana == "X":
        qc.h(0)              # Ana mide en {|+>, |->}
    if base_bruno == "X":
        qc.h(1)              # Bruno mide en {|+>, |->}
    # Ana y Bruno actúan sobre qubits distintos: medir al
    # final da la misma estadística que medir primero a Ana
    # y después a Bruno.
    qc.measure([0, 1], [0, 1])
    return qc

for ana in ["Z", "X"]:
    for bruno in ["Z", "X"]:
        qc = circuito(ana, bruno)
        res = sampler.run([qc], shots=SHOTS).result()[0]
        counts = res.data.c.get_counts()
        # Clave "b1b0": b1 = resultado de Bruno, b0 = de Ana
        conj = {a + b: 0 for a in "01" for b in "01"}
        for k, v in counts.items():
            conj[k[1] + k[0]] = v     # orden (Ana, Bruno)
        bruno_0 = conj["00"] + conj["10"]
        bruno_1 = conj["01"] + conj["11"]
        print(ana, bruno, bruno_0, bruno_1, conj)
\end{verbatim}

Con 4000 disparos y semilla fija, se obtuvieron los conteos de la tabla siguiente. Las columnas de la derecha dan los conteos conjuntos, en el orden (resultado de Ana, resultado de Bruno).
\begin{center}
\renewcommand{\arraystretch}{1.3}
\begin{tabular}{cc|cc|cccc}
  Base de Ana & Base de Bruno & Bruno ve $0$ & Bruno ve $1$ & $00$ & $01$ & $10$ & $11$\\\hline
  $\{\ket0,\ket1\}$ & $\{\ket0,\ket1\}$ & $1999$ & $2001$ & $1999$ & $0$ & $0$ & $2001$\\
  $\{\ket0,\ket1\}$ & $\{\ket+,\ket-\}$ & $1972$ & $2028$ & $1002$ & $1009$ & $970$ & $1019$\\
  $\{\ket+,\ket-\}$ & $\{\ket0,\ket1\}$ & $1973$ & $2027$ & $968$ & $994$ & $1005$ & $1033$\\
  $\{\ket+,\ket-\}$ & $\{\ket+,\ket-\}$ & $2032$ & $1968$ & $2032$ & $0$ & $0$ & $1968$
\end{tabular}
\end{center}
Bruno observa frecuencias cercanas a $\frac12$ en los cuatro casos, con diferencias del orden de $\sqrt{4000\cdot\frac14}\approx32$ disparos. Las frecuencias conjuntas, en cambio, muestran que los resultados de Ana y Bruno están perfectamente correlacionados cuando ambos miden en la misma base, y que no lo están en el caso de bases distintas.
